% =====================================================================
\section{Time scales}
\figp[0.9\textwidth]{fig_times.png}{Characteristic times of the phenomena discussed in the brochure. The scale is logarithmic.}
The processes relevant to the illness span more than ten orders of magnitude in time, from the millisecond impulse of a neuron to the multi-year maturation of the cortex (fig.~2). To the slow processes, the fast ones look instantaneously equilibrated; to the fast ones, the slow look like constant parameters.
\facts
\begin{itemize}
\item A neuronal impulse lasts about a millisecond.\tU\ A hallucination is usually seconds; only 7.4\,\% of episodes last 10 minutes or longer.\tS
\item The amount of a protein in a cell changes over hours and days.\tU
\item The brain's capacity to synthesise dopamine rises over months and years, from prodrome to psychosis.\tS
\item External rhythms impose their own scales: the day (light, sleep and wakefulness) and the year (day length, temperature).\tU
\end{itemize}
\begin{ineq}
A separation of time scales allows the model to be built in layers: fast variables are eliminated as algebraic relations, slow ones are frozen as parameters. One or two slow variables will not cover a spread this wide; the project's model describes the middle of the scale --- minutes to days, events inside a single cell.
\end{ineq}
